% One-column working draft for a proposed JKPS Research Article.
% Journal template conversion and scientific tasks B01--B12 remain pending.
\documentclass[11pt,a4paper]{article}
\usepackage[T1]{fontenc}
\usepackage[utf8]{inputenc}
\usepackage[margin=25mm]{geometry}
\usepackage{amsmath,amssymb}
\usepackage{booktabs,tabularx}
\usepackage[nopatch=footnote]{microtype}
\usepackage{xurl}
\usepackage[hidelinks]{hyperref}
\usepackage{setspace}
\setstretch{1.1}
\newcommand{\draftnote}[1]{\par\smallskip\noindent\textit{Draft revision: #1}\par\smallskip}
% Publication values exported from archived BTS artifacts on 2026-10-13.
% See draft_provenance.json; no new simulations were run.
\newcommand{\BtsPeakX}{31.4368}
\newcommand{\BtsPeakY}{50.7459}
\newcommand{\BtsMedianX}{0.00208629}
\newcommand{\BtsMedianXLo}{0.00207403}
\newcommand{\BtsMedianXHi}{0.00209890}
\newcommand{\BtsDispResidual}{4.382}
\newcommand{\BtsMerit}{1.2906}
\newcommand{\BtsStrengthRows}{%
q11 & +0.448572 & -0.367525 & [-2.5, 2.5] \\
q12 & -1.026778 & -1.216433 & [-3.0, 3.0] \\
q13 & +0.887640 & +1.586072 & [-3.0, 3.0] \\
q21 & -1.066465 & -1.068373 & [-3.0, 3.0] \\
q22 & +1.488384 & +2.823108 & [-3.0, 3.0] \\
q23 & -0.669894 & -2.490968 & [-3.0, 3.0] \\
q31 & +0.589886 & +2.042678 & [-3.0, 3.0] \\
q32 & -1.168702 & -1.434271 & [-3.0, 3.0] \\
q33 & +0.941655 & +0.351115 & [-2.5, 2.5] \\
}
\newcommand{\BtsQuantileRows}{%
$\mathcal M_x$ [1] & 0.002086 & [0.002074, 0.002099] & 0.016482 & 0.028814 \\
$\mathcal M_y$ [1] & 0.002499 & [0.002483, 0.002514] & 0.020701 & 0.035932 \\
$\beta_{x,\max}$ [m] & 31.438418 & [31.430758, 31.445751] & 35.564889 & 37.521376 \\
$\beta_{y,\max}$ [m] & 50.751986 & [50.746458, 50.757604] & 53.831455 & 55.275722 \\
}
\newcommand{\BtsSeedRows}{%
42 & 50,000 & 0 & 0.00210273 & $20.830\times10^{-6}$ & Not met \\
123 & 50,000 & 1 & 0.00207934 & $11.436\times10^{-6}$ & Not met \\
777 & 50,000 & 0 & 0.00210260 & $10.567\times10^{-6}$ & Not met \\
2026 & 400,000 & 4 & 0.00208351 & $2.395\times10^{-6}$ & Met \\
}

\title{Constrained optical matching and error sensitivity of a booster-to-storage-ring transfer line}
\author{Chong Shik Park\\[3pt]
\small Department of Accelerator Science and Center for Accelerator Research,\\
\small Korea University, Sejong, 30019 Republic of Korea\\
\small \texttt{kuphy@korea.ac.kr}}
\date{October 13, 2026}
\begin{document}
\maketitle
\noindent\textbf{Working draft for Journal of the Korean Physical Society.}
This version reports archived simulation results. Italic revision notes identify
analyses and author inputs still needed; they are not completed findings.

\begin{abstract}
Optical matching in a booster-to-storage-ring transfer line must deliver a
specified phase-space distribution while controlling the beam envelope. We
examine a 4 GeV transfer-line model with nine independently adjustable
quadrupoles, using normalized optical residuals and bounded numerical searches
followed by physical feasibility checks. For the archived exit target, the
selected solution matches both transverse planes to floating-point precision,
with peak horizontal and vertical beta functions of 31.44 m and 50.75 m,
respectively, below the imposed 60 m ceiling. Four fixed-optics error ensembles
containing 550,000 realizations give a median horizontal mismatch of 0.002086,
with a 95\% interval of 0.002074--0.002099. Five realizations violate the implemented
optical feasibility criteria, corresponding to a probability estimate of
$9.09\times10^{-6}$ with an exact 95\% interval of
$2.95\times10^{-6}$--$2.12\times10^{-5}$. The historical sensitivity ranking places
input beta errors ahead of quadrupole gradient and energy errors under the
assumed budget. These results characterize an implemented optics model rather
than physical beam loss or storage-ring capture. They establish a matched
reference and quantify statistical precision, while identifying the boundary,
error-coverage and transport validation needed to turn the model into practical
operating tolerances.
\end{abstract}
\noindent\textbf{Keywords:} transfer-line optics; quadrupole matching; beam transport;
error sensitivity; accelerator simulation

\section{Introduction}
Low-emittance storage rings require careful control of the beam distribution
presented to their injection region. The design of diffraction-limited rings
involves strong focusing and limited nonlinear acceptance~\cite{Hettel2014}.
The Korea Light Source design provides a relevant 4 GeV setting for this
problem~\cite{Kim2025}; booster injection and ramping have also been studied for
the Korea-4GSR project~\cite{Lee2024}. A booster-to-storage-ring (BTS) line must
transport the extracted beam to a declared boundary with suitable transverse
optics and dispersion. Its performance depends on both the nominal match and
the response to errors in the incoming beam and magnets.

A small exit mismatch does not alone establish a useful transfer-line design.
Large beta functions elsewhere in the line can reduce aperture margin, while
sensitive quadrupole combinations can make a nominal solution difficult to
maintain. Moreover, the significance of any reported tolerance depends on which
errors actually act in the evaluator and which observables define failure.
Matching, transmission through the BTS and capture in the storage ring are
separate quantities.

This draft examines the deterministic matching and fixed-optics tolerance
campaign archived in the nkm-injection repository. The central question is
whether its nine-quadrupole model reaches a specified exit target within the
configured limits, and what the existing error ensemble establishes about that
solution. We report the selected configuration, its statistical uncertainty and
the limits of the implemented model. Nonlinear-kicker field effects, stored-beam
transparency and multiturn storage-ring capture belong to the companion NKM
study and are outside the present results.

Published transfer-line work emphasizes physical observables and diagnostics.
The SLS 2.0 line separates beam characterization from trajectory control~\cite{Kallestrup2025};
ESRF--EBS rematching includes off-axis injection-quadrupole fields and mechanical
constraints~\cite{ESRF2018}. Sirius studies distinguish transfer efficiency from
ring injection and incoming beam-size effects~\cite{Resende2022}. Pulsed-multipole
injection at MAX IV concerns the downstream capture scheme~\cite{Leemann2012},
which is outside this article. These studies do not supply the boundary values
of the present model.

Our physical question is whether an alternative quadrupole setting, matched to
the same parametric exit target, increases five-sigma clearance at the declared
aperture stations when dispersion and centroid errors are included. We evaluate
a frozen design on held-out static machines and independent shot jitter. This
conditional design comparison makes no claim of algorithm novelty, operational
feedback or measured hardware tolerances.

\section{Transfer-line model and boundary conditions}
\subsection{Lattice and coordinates}
The model is constructed with Accelerator Toolbox and its Python interface,
pyAT~\cite{Terebilo2001,Rogers2017}. Its nominal energy is 4 GeV and its length
is 21.789 m. Nine 0.2 m quadrupoles are arranged in three triplets, named
q11--q13, q21--q23 and q31--q33. The sequence also contains an extraction kicker,
two septum elements and three bending magnets. The nominal bending angles in
sequence are $+7.500$, $+88.500$, $-111.701$, $+176.000$, $-111.701$ and
$+88.500$ mrad; magnet lengths and edge angles are part of the archived
configuration.

Positions are expressed in metres and normalized quadrupole strengths
$K=G/(B\rho)$ in $\mathrm{m}^{-2}$, where $G$ is the magnetic gradient in T/m and
$B\rho$ is the magnetic rigidity in T\,m. Particle tracking uses AT coordinates
$(x,p_x,y,p_y,\delta,ct)$, where $p_x$ and $p_y$ are normalized canonical momenta
and $\delta$ is the relative momentum deviation. Linear transverse optical
quantities use the corresponding paraxial convention. The dispersion derivative
$D'_x$ is dimensionless.

The line configuration places square aperture elements with half-widths 19.35 mm
at its ends and 30 mm at two interior locations. These are model apertures;
continuous vacuum-chamber geometry and measured clearances have not been
established for the present article. The preliminary MAD-X file has ten
quadrupoles and cannot serve as an independent equivalent-lattice validation
without reconciling its sequence.

\subsection{Entrance and exit optics}
Table~\ref{tab:boundary} gives the boundary values used in the archived matching
configuration. The target is treated here as a declared BTS-exit target. It
must not be identified with the storage-ring optics at the NKM centre, since
those values and the intervening transport differ.

\begin{table}[!htbp]
\centering
\caption{Archived entrance optics and declared exit target. Beta and dispersion
are in m; alpha and dispersion derivative are dimensionless. The physical
provenance of the target boundary remains to be confirmed.}
\label{tab:boundary}
\begin{tabular}{lrr}
\toprule Quantity & Entrance & Exit target\\\midrule
$\beta_x$ [m] & 7.560000 & 2.336495\\
$\alpha_x$ [1] & 1.523100 & $-0.016335$\\
$\beta_y$ [m] & 12.269000 & 4.256241\\
$\alpha_y$ [1] & $-1.654700$ & 0.017772\\
$D_x$ [m] & 0.276200 & 0.080868\\
$D'_x$ [1] & $-0.065700$ & 0.047472\\
\bottomrule
\end{tabular}
\end{table}

The historical constraint configuration assumes horizontal and vertical
emittances of 100 and 10 nm\,rad and a relative energy spread of 0.0011. These
are sensitivity-study inputs. Later author documentation identifies a
7.67 nm\,rad horizontal design booster emittance; the archived ensembles are
not relabelled as simulations of that beam. A design-beam transmission study
requires a consistent specification of emittance, coupling, momentum spread,
centroid, distribution and extraction optics.

\draftnote{B02--B04: confirm the source and observation point of the boundary
optics, reconcile AT and MAD-X geometry, provide a sourced layout figure and
validate transport independently. Establish the main beam specification and
physical chamber before reporting design-beam transmission.}

\section{Matching and uncertainty methods}
\subsection{Normalized matching and feasibility selection}
For a quadrupole vector $\mathbf K$ and exit observables
$\mathbf O=(\beta_x,\beta_y,\alpha_x,\alpha_y,D_x,D'_x)$, the implemented objective is
\begin{equation}
 r_i(\mathbf K)=\frac{O_i(\mathbf K)-O_{i,\mathrm{target}}}{s_i},
 \qquad J(\mathbf K)=\sum_{i=1}^{6}r_i^2.
 \label{eq:merit}
\end{equation}
The normalization scales are 0.05 m for each beta residual, 0.01 for each alpha
residual, 0.002 m for dispersion and 0.001 for its derivative. These make
$J$ dimensionless; they are numerical scales rather than demonstrated operating
tolerances.

The archived workflow first uses bounded least squares. When its result is
successful and away from the strength bounds, SLSQP refines the same scalar
merit. The configured budget is 100 for the least-squares maximum function
evaluations and SLSQP maximum iterations. The current SLSQP call supplies
strength bounds but no nonlinear constraint functions. Other physical criteria
are evaluated separately, and the final selection prioritizes feasible
candidates. This distinction matters when comparing the procedure with a
fully constrained search.

The outer random seed is 42. The first of three starts uses nominal settings;
two additional starts are sampled within the family bounds. q11 and q33 have
bounds $[-2.5,2.5]$ m$^{-2}$, while the other families use $[-3,3]$ m$^{-2}$.
The archived constraint configuration also specifies a 60 m beta ceiling,
1.5 m dispersion ceiling, a 1.2 T pole-tip field limit and beam-envelope and
injection-geometry criteria. Their interpretation as hardware specifications
requires source confirmation. Successful optimizer termination and physical
feasibility are recorded separately.

For uncoupled Twiss comparison, let
\begin{equation}
 T_u=\begin{pmatrix}\beta_u&-\alpha_u\\-\alpha_u&
 (1+\alpha_u^2)/\beta_u\end{pmatrix},\qquad
 \mathcal M_u=\tfrac12\operatorname{Tr}
 (T_{u,\mathrm{target}}^{-1}T_{u,\mathrm{out}})-1,
 \label{eq:mismatch}
\end{equation}
where $u=x,y$. The matrices are normalized to equal emittance;
$\mathcal M_u$ is dimensionless and vanishes for a matched plane.

\subsection{Archived tolerance ensemble}
Three independent runs use 50,000 realizations each with seeds 42, 123 and 777;
a follow-up uses 400,000 realizations with seed 2026. All select the same
optimization and error-budget hashes. The quadrupole settings remain fixed
while errors are sampled. The analysis contains no capture callback and no
particle-through-aperture transmission measurement. Its optical failure
criterion is
\begin{equation}
 \max(\beta_{x,\max},\beta_{y,\max})>60.01~\mathrm m
 \quad\text{or}\quad \mathcal M_x+\mathcal M_y>2.05,
 \label{eq:failure}
\end{equation}
including the stored 0.01 m beta and 0.05 mismatch allowances. This definition
is less extensive than the nominal design's full feasibility checks.

Quantile intervals are calculated from order statistics using a binomial
coverage construction. Median bootstrap checks use 20,000 replicates, with
resampling seeds 1729, 2718 and 31415. For $k$ failures in $N$ independent
realizations, the two-sided exact 95\% interval is obtained from beta-distribution
quantiles at 0.025 and 0.975 with parameters $(k,N-k+1)$ and $(k+1,N-k)$,
respectively, with endpoint conventions for $k=0$ or $k=N$. Wilson score
intervals provide a cross-check~\cite{Wilson1927}. These intervals describe
sampling uncertainty conditional on the chosen model and budget, not uncertainty
in the physical assumptions.

\subsection{Error coverage}
The historical budget includes relative quadrupole-gradient errors with
standard deviation $10^{-3}$, relative energy errors of $10^{-3}$, input beta
errors of 5\%, quadrupole offsets of 100 $\mu$m and rolls of 0.5 mrad.
The coverage audit reports effects on the optics from gradient, energy, beta
and roll terms, with an offset contribution to peak beta through the implemented
approximation. Dipole-strength errors, input centroid jitter and longitudinal
placement errors do not affect the reported optical metrics in that evaluator.
The same driver also computes an NKM stored-kick diagnostic; this diagnostic is
excluded from the BTS results.

Consequently, an inactive error cannot be assigned a generous physical tolerance
from a zero numerical response. Coupled transport and orbit-dependent clearance
need explicit validation, particularly for roll and alignment errors.
\draftnote{B07: generate a machine-readable perturbation coverage audit and
implement the missing physical responses needed for BTS transport claims.
Document distributions, correlations and sources for each error amplitude.}

\section{Matched design and archived performance}
\subsection{Selected quadrupole settings}
The archived search finds one feasible candidate among three starts, with the
termination message ``Optimization terminated successfully.'' Its final merit
is $\BtsMerit\times10^{-13}$, and both mismatch values print as zero at
floating-point precision. The horizontal dispersion residual is
$\BtsDispResidual\times10^{-11}$ m. This is a numerical match to the declared
target, not evidence of a global optimum or an experimentally realized setting.
Table~\ref{tab:strengths} reports the selected strengths.

\begin{table}[!htbp]
\centering
\caption{Nominal and selected quadrupole strengths, and family bounds,
all in m$^{-2}$. Values are exported from the selected archived summary and
configuration. Rounded settings have not been independently revalidated here.}
\label{tab:strengths}
\begin{tabular}{lrrc}
\toprule Family & Nominal $K$ & Selected $K$ & Bounds\\\midrule
\BtsStrengthRows
\bottomrule
\end{tabular}
\end{table}

The selected peak beta functions are $\BtsPeakX$ m horizontally and
$\BtsPeakY$ m vertically. Relative to the configured 60 m ceiling, the
vertical plane has the smaller beta margin. That margin is not itself a
physical clearance: for an uncoupled beam, the rms horizontal size obeys
\begin{equation}
 \sigma_x^2(s)=\beta_x(s)\varepsilon_x+D_x^2(s)\sigma_\delta^2.
 \label{eq:envelope}
\end{equation}
A chamber check must also include the centroid orbit, vertical or coupled
covariance, a declared beam-coverage factor and the aperture shape.

\draftnote{B05 and B08: recompute baseline and matched optics from the same
frozen configuration; compare a design chosen for envelope or error margin
using equal physical assumptions. Add optics/dispersion and clearance profiles.
No new robust design or nominal transmission gain has yet been established.}

\subsection{Tolerance distributions}
All 550,000 requested evaluations in the selected historical ensemble are
numerically valid. Table~\ref{tab:quantiles} gives pooled mismatch and peak-beta
quantiles, omitting the stored-kick observable. The median horizontal mismatch
is $\BtsMedianX$, with a 95\% interval
$[\BtsMedianXLo,\BtsMedianXHi]$. Its upper tail is appreciably larger than its
median, so a single central estimate does not describe every realization.

\begin{table}[!htbp]
\centering\small
\caption{Pooled fixed-optics results from 550,000 archived realizations.
The interval is the 95\% order-statistic interval for the median. Mismatch is
dimensionless. These metrics do not measure particle transmission.}
\label{tab:quantiles}
\begin{tabular}{lrrrr}
\toprule Observable & Median & Median interval & p95 & p99\\\midrule
\BtsQuantileRows
\bottomrule
\end{tabular}
\end{table}

Five realizations fail the optical criteria in Eq.~\eqref{eq:failure}, giving
$\hat p=5/550000=9.09\times10^{-6}$ and an exact 95\% interval of
$[2.95\times10^{-6},2.12\times10^{-5}]$. The small count leaves considerable
relative uncertainty in this probability. It must not be interpreted as
five lost bunches, an accelerator reliability specification, or a measurement
of ring injection efficiency.

\section{Error sensitivity and statistical precision}
\subsection{Independent seeds and prefix checks}
Table~\ref{tab:seeds} records the independent runs and their predeclared absolute
median-stability target of $10^{-5}$. The 50,000-sample runs compare prefixes
of 25,000 and 50,000 samples; the follow-up compares 200,000 and 400,000.
All three smaller runs miss the target. The follow-up meets it with an
absolute change of $2.395\times10^{-6}$.

\begin{table}[!htbp]
\centering\small
\caption{Per-run sample counts, optical failure counts, median horizontal
mismatch and prefix changes. ``Met'' refers only to the $10^{-5}$ absolute
median-stability target, not to convergence of rare failures or tail quantiles.}
\label{tab:seeds}
\begin{tabular}{rrrrrc}
\toprule Seed & Samples & Failures & Median $\mathcal M_x$ & Prefix change & Target\\\midrule
\BtsSeedRows
\bottomrule
\end{tabular}
\end{table}

A passing median-prefix check does not certify every statistic. The pooled
median interval has a half-width of about $1.24\times10^{-5}$, while the
failure-probability interval remains broad relative to its estimate. Sampling
precision must therefore be assessed separately for central mismatch, extreme
envelopes and rare feasibility failures. Increasing bootstrap count cannot
replace independent physical realizations or correct an incomplete error model.

\subsection{Sensitivity interpretation and correction analysis}
Under the historical budget, the one-at-a-time ranking places input beta
mismatch first, followed by quadrupole-gradient and energy errors. This ordering
is conditional on the specified error amplitudes and the scalar merit.
It does not rank errors absent from the optical evaluator, nor does it determine
an optimal hardware specification independently of beam and aperture assumptions.

An interpretable correction study should use a dimensionless response matrix,
\begin{equation}
 A_{ij}=\frac{\Delta K_{j,\mathrm{ref}}}{s_i}
 \frac{\partial O_i}{\partial K_j},\qquad A=U\Sigma V^{\mathsf T},
 \label{eq:svd}
\end{equation}
where $s_i$ are declared optical scales and $\Delta K_{j,\mathrm{ref}}$ are
meaningful strength-change scales. Its singular vectors can identify correction
combinations and weakly observable freedoms, but finite-step reevaluation and
constraint checks are needed before treating them as usable knobs. The raw
mixed-unit Jacobian condition number is not used as a tolerance result here.

\draftnote{B06, B09 and B10: validate finite-difference steps and singular rank,
measure finite correction range, scan relevant error amplitudes, and compare
uncorrected/corrected transport on held-out static errors with independent shot
jitter. Add sensitivity and correction figures only from completed analyses.}

\section{Discussion}
The archived solution demonstrates that the chosen linear BTS model can meet
its declared optical target with peak beta below the configured ceiling.
The error ensemble provides a reproducible conditional characterization of
mismatch and envelope fluctuations. Its principal practical message is that
input optics must be controlled alongside magnet settings; a nominally exact
match does not remove sensitivity to the extracted beam.

Several limits prevent direct translation to commissioning requirements.
The target boundary and lattice provenance need source confirmation; the
historical beam emittance differs from the later design specification; and
some sampled errors have no effect on the reported optical quantities.
The nominal feasibility validator and the ensemble optical-failure rule also
test different sets of conditions. Physical transmission requires explicit
particles and apertures, including orbit and coupling effects, rather than a
beta threshold alone.

The archived multi-objective study is not used to support the present results.
Its dispersion objective combines $D_x$ in metres with dimensionless $D'_x$,
and its legacy finalist procedure does not establish physical BTS transmission.
A future trade-off study should normalize both residuals consistently and
compare physically validated designs under common constraints and held-out
errors. It need not use a genetic algorithm if deterministic searches expose
the relevant margin trade-off.

\draftnote{B08--B10: replace these prospective points with a quantified physical
design and correction comparison before claiming the planned article's additional
contribution. B01--B12 remain pending; this first draft does not complete them.}

\section{Conclusions}
For the archived 4 GeV nine-quadrupole BTS model, the selected bounded-search
solution matches the declared exit optics to floating-point precision and has
peak beta functions of 31.44 and 50.75 m. Under the implemented fixed-optics
error budget, 550,000 realizations yield a median horizontal mismatch of
0.002086 and five optical feasibility failures. The larger follow-up meets the
specified median-prefix target, while rare-failure uncertainty remains distinct.
These findings establish a reference for further transport and correction
studies. They do not establish physical beam transmission, hardware tolerances
or storage-ring capture. A source-confirmed boundary model, active-error audit
and held-out margin/correction comparison are needed to substantiate the final
BTS article's engineering conclusions.

\section*{Statements and declarations}
\paragraph{Author contributions.}
\draftnote{Author to confirm contributions and final authorship.}
\paragraph{Funding and competing interests.}
\draftnote{Author to supply funding/grant information and confirm the competing-interest declaration.}
\paragraph{Acknowledgments and generative assistance.}
This working draft was prepared with OpenAI Codex assistance. Scientific claims
are limited to the selected archived artifacts; the author must verify the text
and approve the final assistance disclosure and acknowledgments.
\paragraph{Data and code availability.}
The model and analysis code are maintained in the nkm-injection repository
(\url{https://github.com/cspark7701/nkm-injection}). The selected simulation
outputs are ignored by version control and are not supplied by the repository
alone. A local provenance record accompanies this draft. No permanent public
archive or BTS referee attachment is claimed to exist.
\draftnote{B12: provide the selected raw-output attachment, source/environment
snapshot and independently verified rebuild instructions before submission.}

\appendix
\section{Draft evidence selection and unfinished assets}
The nominal solution is selected from
\path{backups_20261009/production_statistical_01/optimization/}; its summary and
adjacent configuration match the hashes recorded in the outline. The tolerance
selection is
\path{results/journal_campaign_01/followup_analysis/tolerance_replicates_with_400k.json}.
All four raw-summary hashes were checked when preparing the publication values.
The exported tables are installed in \path{generated/archived_results.tex};
\path{draft_provenance.json} records the source paths, digests and source commit.
The current PDF can be compiled from these installed publication files without
access to the ignored raw ensembles, but such compilation does not rerun or
independently validate their physics.

The six planned main figures are not inserted as simulated or measured results.
The layout, optics/envelope profiles, normalized correction responses, margin
trade-off, active-error scans and held-out correction/distribution plots remain
assigned to B02--B10 in \path{002_manuscript_tasks.md}. The bibliography below
contains six entries selected from the existing audited reference list; the
closest BTS literature comparison remains B01 work. This generic one-column
article format is for drafting; final JKPS formatting is a B11 deliverable.

\bibliographystyle{unsrt}
\bibliography{paper}
\end{document}
